% Generated draft: requires mathematical and source review.
\documentclass[11pt]{article}
\usepackage[margin=25mm]{geometry}
\usepackage{fontspec}
\setmainfont{Latin Modern Roman}
\usepackage{amsmath,amssymb,mathtools,unicode-math}
\setmathfont{Latin Modern Math}
\usepackage{enumitem,needspace,xurl,hyperref}
\setlength{\parindent}{0pt}
\setlength{\parskip}{5pt}
\setlength{\emergencystretch}{3em}
\newcommand{\N}{\mathbb{N}}
\newcommand{\Z}{\mathbb{Z}}
\newcommand{\Q}{\mathbb{Q}}
\newcommand{\R}{\mathbb{R}}
\newcommand{\C}{\mathbb{C}}
\newcommand{\F}{\mathbb{F}}
\newcommand{\A}{\mathbb{A}}
\newcommand{\PP}{\mathbb{P}}
\newcommand{\E}{\mathbb{E}}
\newcommand{\eps}{\varepsilon}
\newcommand{\dd}{\,\mathrm{d}}
\begin{document}
\section*{UnsolvedMath ID: 30000267}
\textbf{Source problem code:} catalogue entry
\subsection{Definitions and mathematical statement}
\textbf{Multiplicity conjecture.} Consider the stationary nonlinear Schr\"odinger equation\[-\Delta u+(\lambda^2a(x)+1)u=|u|^{p-1}u\quad\text{on }\R^N,\] with two disjoint zero-potential regions \(\Omega_1,\Omega_2\). For all sufficiently large \(\lambda\), must there be at least\[\operatorname{cat}(\Omega_1)+\operatorname{cat}(\Omega_2)+\operatorname{cat}(\Omega_1)\operatorname{cat}(\Omega_2)\] distinct positive solutions? The large-parameter quantifier means a threshold \(\lambda_0\), depending on the fixed admissible problem data, such that the bound holds for every \(\lambda>\lambda_0\); it is not merely an assertion along one sequence of parameters. Positivity means \(u>0\) in the solution convention selected for this equation.\par Under the unnormalized Lusternik--Schnirelmann convention, \(\operatorname{cat}(\Omega)\) is the least number of open subsets covering \(\Omega\) whose inclusions into \(\Omega\) are null-homotopic; in particular a nonempty contractible space has category one. The source does not specify this convention, and using a normalized category would change the numerical bound. Finiteness of the categories and the regularity and topology of the wells must also be specified.\par The excerpt supplies no range of \(N,p\), sign or regularity hypotheses on \(a\), conditions at infinity, or solution space. In the usual nonnegative-potential finite-energy formulation one would require \(u\in H^1(\R^N)\) and \(\int_{\R^N}a(x)|u|^2\,dx<\infty\), with the PDE satisfied weakly; this is an explanatory possible convention, not a recovered source hypothesis. The relevant superlinear or subcritical exponent restrictions are not imposed silently. The phrase \emph{zero-potential region} could mean \(\{a=0\}\), its interior, or the interiors of components of its closure. In particular, for continuous \(a\), its zero set is closed, so identifying it literally with a union of open wells requires care. The location, boundedness, boundary regularity, and separation of the two regions are missing. These gaps prevent an unconditional theorem for every potential matching only the short description.\par The supplied triage checked 2026-08-21 labels the conjecture open. It reports related category-based multiplicity results but no located proof of this exact bound in the extracted generality. This is a dated search assessment, not independent evidence of the current sharp multiplicity theorem.
\subsection{Short English statement}
For a large-parameter equation with two potential wells, ask for the stated sum-and-product category lower bound on distinct positive solutions. The category normalization, potential and exponent hypotheses, and finite-energy solution convention are not supplied and cannot be chosen silently.
\subsection{Sources}
\begin{itemize}\item Local UnsolvedMath registry, permanent ID 30000267, code catalogue entry; \emph{Category Bounds for Positive Nonlinear Schr\"odinger Solutions}, Oberwolfach Reports Open Problems; original label \(P_\lambda\).\item \emph{Topological and Variational Methods for Differential Equations (2006)}, Oberwolfach Reports, DOI 10.4171/owr/2005/29, pp. 1601--1678. \url{https://doi.org/10.4171/owr/2005/29}. Proposed year and year tag: 2005; no proposer named.\item Supplied literature label \emph{original problem report}: \url{https://doi.org/10.4171/owr/2005/29}.\item Literature triage checked 2026-08-21, status \texttt{open}; record updated 2026-09-16. The \texttt{corrected\_verified} field reports restoring the equation and two-well hypothesis from the official report, not independently checked here.\end{itemize}
\end{document}
